% !TEX program = xelatex
\documentclass[11pt,a4paper]{ctexart}
\usepackage[a4paper,margin=2.15cm]{geometry}
\usepackage{fontspec}
\usepackage{amsmath,amssymb}
\usepackage{xcolor}
\usepackage{booktabs,tabularx,longtable,array,multirow}
\usepackage{enumitem}
\usepackage{fvextra}
\usepackage{newunicodechar}
\usepackage[most]{tcolorbox}
\usepackage{url}
\usepackage[colorlinks=true,linkcolor=blue!50!black,urlcolor=blue!55!black,citecolor=blue!50!black]{hyperref}
\usepackage{fancyhdr}

% ---------- fonts: monospace with wide Unicode coverage for Lean ----------
\setmonofont{DejaVuSansMono.ttf}[Scale=0.80,BoldFont=DejaVuSansMono-Bold.ttf,ItalicFont=DejaVuSansMono-Oblique.ttf,BoldItalicFont=DejaVuSansMono-BoldOblique.ttf]
\newfontfamily\mathfb{STIX Two Math}[Scale=0.88]
\newunicodechar{𝒞}{{\mathfb 𝒞}}
\newunicodechar{𝓞}{{\mathfb 𝓞}}
\newunicodechar{𝒜}{{\mathfb 𝒜}}
\newunicodechar{𝒮}{{\mathfb 𝒮}}

% ---------- macros ----------
\makeatletter\g@addto@macro{\UrlBreaks}{\do\_\do\.}\makeatother
\Urlmuskip=0mu plus 2mu
\DeclareUrlCommand\lean{\urlstyle{tt}}
\DeclareUrlCommand\path{\urlstyle{tt}}
\newcommand{\leanu}[1]{\texttt{#1}}
\newcommand{\mtt}[1]{\text{\texttt{#1}}}% identifiers with non-ASCII characters
\newcommand{\OK}{\textcolor{green!45!black}{\textbf{属实}}}
\newcommand{\PARTLY}{\textcolor{orange!70!black}{\textbf{部分属实}}}
\newcommand{\NO}{\textcolor{red!70!black}{\textbf{不成立}}}
\newcommand{\NOTFOLLOW}{\textcolor{red!70!black}{推不出}}
\definecolor{boxgray}{gray}{0.955}
\newtcolorbox{verdictbox}[1][]{enhanced,breakable,colback=blue!3,colframe=blue!45!black,boxrule=0.8pt,arc=2pt,left=7pt,right=7pt,top=6pt,bottom=6pt,#1}
\newtcolorbox{notebox}[1][]{enhanced,breakable,colback=boxgray,colframe=gray!55,boxrule=0.5pt,arc=2pt,left=7pt,right=7pt,top=5pt,bottom=5pt,#1}
\newtcolorbox{warnbox}[1][]{enhanced,breakable,colback=orange!6,colframe=orange!70!black,boxrule=0.6pt,arc=2pt,left=7pt,right=7pt,top=5pt,bottom=5pt,#1}

\makeatletter\newcommand{\tableinput}[1]{\@@input #1 }\makeatother
\expandafter\def\csname ex@anygeo@first\endcsname{14}
\expandafter\def\csname ex@anygeo@last\endcsname{19}
\expandafter\def\csname ex@anygeo@file\endcsname{\path{src/Development/AbelianVarietyTheory/Definitions/AutomorphicMembership/OverIx.lean}}
\expandafter\def\csname ex@oat@first\endcsname{44}
\expandafter\def\csname ex@oat@last\endcsname{49}
\expandafter\def\csname ex@oat@file\endcsname{\path{src/Development/AbelianVarietyTheory/Definitions/AutomorphicMembership/OverIx.lean}}
\expandafter\def\csname ex@geocl@first\endcsname{21}
\expandafter\def\csname ex@geocl@last\endcsname{28}
\expandafter\def\csname ex@geocl@file\endcsname{\path{src/Development/AbelianVarietyTheory/Definitions/AutomorphicMembership/NotGeoClFalse.lean}}
\expandafter\def\csname ex@main@first\endcsname{1}
\expandafter\def\csname ex@main@last\endcsname{12}
\expandafter\def\csname ex@main@file\endcsname{\path{src/Main.lean}}
\expandafter\def\csname ex@maintheorems@first\endcsname{7}
\expandafter\def\csname ex@maintheorems@last\endcsname{9}
\expandafter\def\csname ex@maintheorems@file\endcsname{\path{src/Supplements/Kernels/MainTheorems.lean}}
\expandafter\def\csname ex@mainconcl@first\endcsname{7}
\expandafter\def\csname ex@mainconcl@last\endcsname{11}
\expandafter\def\csname ex@mainconcl@file\endcsname{\path{src/Development/AutomorphicArithmetic/CharacterBoxes/MainConclusion.lean}}
\expandafter\def\csname ex@lit810@first\endcsname{810}
\expandafter\def\csname ex@lit810@last\endcsname{811}
\expandafter\def\csname ex@lit810@file\endcsname{\path{src/Development/AutomorphicArithmetic/ThetaCorrespondence/ThetaPairLifts/ThirdPlaceSeries/EveryCharacterOfLiteratureOfVWin.lean}}
\expandafter\def\csname ex@siegel_a@first\endcsname{59}
\expandafter\def\csname ex@siegel_a@last\endcsname{67}
\expandafter\def\csname ex@siegel_a@file\endcsname{\path{src/Development/AutomorphicArithmetic/ThetaCorrespondence/ThetaPairLifts/ThirdPlaceSeries/SiegelEisensteinSecondTermHypotheses.lean}}
\expandafter\def\csname ex@siegel_b@first\endcsname{69}
\expandafter\def\csname ex@siegel_b@last\endcsname{73}
\expandafter\def\csname ex@siegel_b@file\endcsname{\path{src/Development/AutomorphicArithmetic/ThetaCorrespondence/ThetaPairLifts/ThirdPlaceSeries/SiegelEisensteinSecondTermHypotheses.lean}}
\expandafter\def\csname ex@topplug@first\endcsname{16}
\expandafter\def\csname ex@topplug@last\endcsname{18}
\expandafter\def\csname ex@topplug@file\endcsname{\path{src/Development/AutomorphicArithmetic/ThetaCorrespondence/ThetaPairLifts/ThirdPlaceSeries/TopPlug.lean}}
\expandafter\def\csname ex@plugall@first\endcsname{16}
\expandafter\def\csname ex@plugall@last\endcsname{18}
\expandafter\def\csname ex@plugall@file\endcsname{\path{src/Development/AutomorphicArithmetic/CharacterBoxes/EveryCharacter.lean}}
\expandafter\def\csname ex@everyE2@first\endcsname{67}
\expandafter\def\csname ex@everyE2@last\endcsname{71}
\expandafter\def\csname ex@everyE2@file\endcsname{\path{src/Development/AutomorphicArithmetic/ThetaCorrespondence/FockModels/EveryCharacterE.lean}}
\expandafter\def\csname ex@prm1@first\endcsname{20}
\expandafter\def\csname ex@prm1@last\endcsname{21}
\expandafter\def\csname ex@prm1@file\endcsname{\path{src/Development/AutomorphicArithmetic/ThetaCorrespondence/ThetaPairLifts/ThirdPlaceSeries/Prm1.lean}}
\expandafter\def\csname ex@params@first\endcsname{19}
\expandafter\def\csname ex@params@last\endcsname{25}
\expandafter\def\csname ex@params@file\endcsname{\path{src/Development/AutomorphicArithmetic/ThetaCorrespondence/ThetaPairLifts/ThirdPlaceSeries/CoreAltAltBis.lean}}
\expandafter\def\csname ex@mono_sig@first\endcsname{28}
\expandafter\def\csname ex@mono_sig@last\endcsname{28}
\expandafter\def\csname ex@mono_sig@file\endcsname{\path{src/Development/AutomorphicArithmetic/ThetaCorrespondence/ThetaPairLifts/ThirdPlaceSeries/OpenAndLinkedOfStrongerExtension.lean}}
\expandafter\def\csname ex@geo0_sig@first\endcsname{18}
\expandafter\def\csname ex@geo0_sig@last\endcsname{18}
\expandafter\def\csname ex@geo0_sig@file\endcsname{\path{src/Development/AutomorphicArithmetic/ThetaCorrespondence/ThetaPairLifts/ThirdPlaceSeries/RowGeo0AtStatement.lean}}
\expandafter\def\csname ex@openorbit@first\endcsname{66}
\expandafter\def\csname ex@openorbit@last\endcsname{66}
\expandafter\def\csname ex@openorbit@file\endcsname{\path{src/Development/AutomorphicArithmetic/ThetaCorrespondence/ThetaPairLifts/ThirdPlaceSeries/OpenOrbitOfUnramifiedExponents.lean}}
\expandafter\def\csname ex@openstmt@first\endcsname{39}
\expandafter\def\csname ex@openstmt@last\endcsname{51}
\expandafter\def\csname ex@openstmt@file\endcsname{\path{src/Development/AutomorphicArithmetic/ThetaCorrespondence/ThetaPairLifts/ThirdPlaceSeries/GeometricWeakTies.lean}}
\expandafter\def\csname ex@hypsalt@first\endcsname{22}
\expandafter\def\csname ex@hypsalt@last\endcsname{30}
\expandafter\def\csname ex@hypsalt@file\endcsname{\path{src/Development/AutomorphicArithmetic/ThetaCorrespondence/ThetaPairLifts/ThirdPlaceSeries/Hyps.lean}}
\expandafter\def\csname ex@divexprov@first\endcsname{204}
\expandafter\def\csname ex@divexprov@last\endcsname{209}
\expandafter\def\csname ex@divexprov@file\endcsname{\path{src/Development/AbelianVarietyTheory/Proofs/IwasawaDecomposition/TupPinsOfMuIsW.lean}}
\expandafter\def\csname ex@divexgeocl@first\endcsname{39}
\expandafter\def\csname ex@divexgeocl@last\endcsname{49}
\expandafter\def\csname ex@divexgeocl@file\endcsname{\path{src/Development/AbelianVarietyTheory/Proofs/IwasawaDecomposition/PublishedInputsImplyConclusion.lean}}
\expandafter\def\csname ex@oatsetdiv@first\endcsname{106}
\expandafter\def\csname ex@oatsetdiv@last\endcsname{106}
\expandafter\def\csname ex@oatsetdiv@file\endcsname{\path{src/Development/AbelianVarietyTheory/Proofs/IwasawaDecomposition/DivisionGeometryScalarsTransferSetDivision.lean}}
\expandafter\def\csname ex@geodegen@first\endcsname{241}
\expandafter\def\csname ex@geodegen@last\endcsname{260}
\expandafter\def\csname ex@geodegen@file\endcsname{\path{src/Formalism/LocalToGlobal/LatticeGeometry/Inhabited.lean}}
\expandafter\def\csname ex@divxfl@first\endcsname{829}
\expandafter\def\csname ex@divxfl@last\endcsname{834}
\expandafter\def\csname ex@divxfl@file\endcsname{\path{src/Formalism/LocalToGlobal/LatticeGeometry/Inhabited.lean}}
\expandafter\def\csname ex@geoover@first\endcsname{155}
\expandafter\def\csname ex@geoover@last\endcsname{155}
\expandafter\def\csname ex@geoover@file\endcsname{\path{src/Development/AbelianVarietyTheory/Proofs/IwasawaDecomposition/TupPinsOfMuIsW.lean}}
\expandafter\def\csname ex@geoover2@first\endcsname{164}
\expandafter\def\csname ex@geoover2@last\endcsname{166}
\expandafter\def\csname ex@geoover2@file\endcsname{\path{src/Development/AbelianVarietyTheory/Proofs/IwasawaDecomposition/TupPinsOfMuIsW.lean}}
\expandafter\def\csname ex@heighteq@first\endcsname{379}
\expandafter\def\csname ex@heighteq@last\endcsname{400}
\expandafter\def\csname ex@heighteq@file\endcsname{\path{src/Development/AbelianVarietyTheory/Definitions/GlobalTwisting/SetUnipotentTranslate.lean}}
\expandafter\def\csname ex@buildaid@first\endcsname{1}
\expandafter\def\csname ex@buildaid@last\endcsname{4}
\expandafter\def\csname ex@buildaid@file\endcsname{\path{src/BuildAid/D3.lean}}
\expandafter\def\csname ex@cptcopy@first\endcsname{68}
\expandafter\def\csname ex@cptcopy@last\endcsname{84}
\expandafter\def\csname ex@cptcopy@file\endcsname{\path{src/Formalism/HermitianCurvature/HermitianCurvatureProofs/CutCptCopy.lean}}
\expandafter\def\csname ex@kcopy@first\endcsname{31}
\expandafter\def\csname ex@kcopy@last\endcsname{42}
\expandafter\def\csname ex@kcopy@file\endcsname{\path{src/Formalism/LocalToGlobal/NumberFieldFamilies/OwnStruct.lean}}

\newcommand{\excerpt}[2][]{%
  \par\smallskip\noindent{\scriptsize\color{black!62}\csname ex@#2@file\endcsname\ \ 第 \csname ex@#2@first\endcsname--\csname ex@#2@last\endcsname\ 行}\par\vspace{-0.55em}%
  \VerbatimInput[numbers=left,firstnumber=\csname ex@#2@first\endcsname,breaklines=true,breakanywhere=true,fontsize=\scriptsize,frame=single,rulecolor=\color{gray!50},numbersep=4pt,xleftmargin=2.1em,framesep=3pt,#1]{excerpts/#2.txt}\vspace{-0.3em}}
\fvset{breaklines=true,breakanywhere=true}
% Filled in from the actual runs (see evidence/).
\newcommand{\conlecheSummary}{我们还在本机用独立检查器 con-leche 重放了完整的导出证明：1,468,418 个声明全部被接受，五个主定理全部被判定为 ``justified by the solution''，
而用于比对的挑战导出是我们自己从官方 Mathlib 重建并逐字节核对过的。}
\newcommand{\conlecheDetail}{%
\begin{verdictbox}
\textbf{结果：通过。}退出码 0。标准输出全文如下：
\VerbatimInput[fontsize=\scriptsize,breaklines=true,breakanywhere=true]{../evidence/con-leche/compare.stdout.log}
\end{verdictbox}

\noindent 即：解答导出的 1,468,418 个声明全部在 \lean{--verified} 模式下被接受；挑战导出的 51,692 个声明被接受，
其中 5 个以 \lean{sorry} 占位的定理——包括 \lean{riemannHypothesis}——全部由解答「证成」。这两个数字与 MIT 页面公布的 con-leche 结果完全一致。
本机运行耗时 73 分 9 秒（解析 340.3 秒，检查 4022.3 秒，16 个工作线程），峰值内存 47.4\,GB；
单个声明中最慢的是 \lean{GlobalTheory.HeightCounting.SideCopy.LamCopy.lam_copy}（374 秒）。日志见 \path{evidence/con-leche/}。

这一步同时回答了两个问题：其一，证明项确实通过了一个独立于官方内核实现的检查器（因而 \S\ref{sec:hygiene} 中 \lean{kexactB}／\lean{kcopy} 所依赖的「内核把关」是经得起独立复核的）；
其二，被证明的五个命题与「只导入 Mathlib 的五个陈述」是同一组命题——连同它们读到的每一个定义。
}
\newcommand{\scanThreeSummary}{第三遍扫描（结构体字段读取与消去）的要点：\lean{OAt.rec}、\lean{DivT.rec}、\lean{DivX.rec}、\lean{DV.Geo.rec} 均为零引用，
即这些记录的内容只可能经字段投影读出；\lean{OAt} 的 \lean{geo} 字段没有转手以外的读取者；内容提取引理 \lean{GeoTiedP.geoTied}、\lean{GeoTiedP.formsTieX} 不在依赖锥内；
而抽象接口 \lean{DV.Geo} 的字段确被锥内 621 个声明读取（\S\ref{sec:fact-degenerate}）。}
\newcommand{\scanThreeAppendix}{%
\subsection{第三遍扫描：结构体字段的读取与消去}

「依赖数」是传递依赖该模式的声明个数；投影类种子匹配的是内核项中的 \lean{proj} 结点。

\begin{center}\small
\begin{tabular}{@{}lrc@{}}
\toprule
种子（模式） & 依赖它的声明数 & 主定理依赖它？\\
\midrule
\tableinput{tables/scan3_counts.tex}
\bottomrule
\end{tabular}
\end{center}

直接读取 \lean{OAt} 的 \lean{lam} 字段的定理有三个：\lean{divEx_geoCl}、\lean{oAt_setDiv}（二者是转手）以及 \lean{GlobalTheory.T.parameter_of_discrete_series_ray_alt}（真正的消费者）；
直接读取 \lean{geo} 字段的只有前两个。实际出现在依赖锥内的 \lean{DV.Geo} 字段访问函数共 18 个：
\lean{Var}、\lean{Cl}、\lean{Dv}、\lean{dvGrp}、\lean{Fn}、\lean{fnField}、\lean{fnAlg}、\lean{H0}、\lean{supp}、\lean{dim}、\lean{Adl}、\lean{ofW}、\lean{inter}、\lean{twist}、\lean{vanishOn}、\lean{arTop}、\lean{arVol}、\lean{nefA}。
}



\setlist{itemsep=0.15em,topsep=0.3em,parsep=0.1em}
\setlength{\emergencystretch}{3.5em}
\pagestyle{fancy}\fancyhf{}
\fancyhead[L]{\small anyGeometry 质疑的独立审计报告}\fancyhead[R]{\small 2026-11-23}
\fancyfoot[C]{\small\thepage}
\renewcommand{\headrulewidth}{0.3pt}
\setlength{\headheight}{14pt}
\hypersetup{pdftitle={anyGeometry := True 能否推翻黎曼猜想的 Lean 证明？独立审计报告},pdfauthor={Claude (Anthropic), commissioned by ArtificialZeng}}

\begin{document}
\begin{titlepage}
\centering
\vspace*{2.2cm}
{\LARGE\bfseries \texttt{anyGeometry := True} 能否推翻\\[0.35em]黎曼猜想的 Lean 证明？\par}
\vspace{0.9cm}
{\Large 对 Jenny Lorraine Nielsen「证明无效」论证的独立审计报告\par}
\vspace{0.5cm}
{\large An independent audit of the ``vacuous geometry predicate'' objection\\ to Alpöge's Lean proof of the Riemann Hypothesis\par}
\vspace{1.6cm}
\begin{tabularx}{0.94\textwidth}{@{}r>{\raggedright\arraybackslash}X@{}}
审计对象： & Levent Alpöge（与 Claude 合作）2026-11-20 发布的 Lean 4 证明库（约 1740 万行，175,187 个模块）\\[0.25em]
被审计的质疑： & J.~L.~Nielsen, \emph{Invalidity of the Alpoge Proposed Proof of the Riemann Hypothesis}\\
 & （PhilArchive, 2026-11-21 起多次更新）\\[0.25em]
审计执行： & Claude（Anthropic，Claude Code 代理），受 ArtificialZeng 委托\\[0.25em]
交叉验证： & 4 个独立的 \texttt{agy} 代理（Gemini 4 Neon, XHigh）\\[0.25em]
审计日期： & 2026 年 11 月 23 日（单次会话，时间预算 7 小时）\\[0.25em]
报告版本： & v1.0
\end{tabularx}
\vfill
\begin{verdictbox}
\textbf{一句话结论：打假不成立。}Nielsen 指出的代码\emph{事实}基本属实——而且我们在内核级证明导出文件上确认了
公开讨论中此前悬而未决的一点：\lean{anyGeometry} 以及代码自称 \emph{degenerate}/\emph{junk} 的几何模型，\emph{确实}位于
\lean{riemannHypothesis} 的依赖锥之内。但由这些事实推出「证明无效／空洞」的\emph{推理}不成立：
最终定理是不带任何前提的闭命题，其类型就是 Mathlib 的 \lean{RiemannHypothesis}；该「几何槽」在整条证明路径上只写不读，
是一个死参数而不是一个漏洞。我们未发现该证明存在任何真正的 reward hacking。
\end{verdictbox}
\vspace{0.6cm}
{\small 本报告的全部证据、脚本与日志随报告一并提交于\\ \url{https://github.com/ArtificialZeng/Riemann-hypothesis-Auditing}\par}
\end{titlepage}

\tableofcontents
\newpage

\section{结论摘要}

\subsection{审计问题}

网络流传的质疑（见委托人提供的截图与新闻）可以归结为一句话：证明库里有这样一行定义
\begin{Verbatim}[fontsize=\small,frame=single,rulecolor=\color{gray!50}]
def anyGeometry : GeoP := fun _ _ _ => True
\end{Verbatim}
\noindent 于是「原本需要认真证明的条件被设置成了自动成立，整个程序因此顺利通过检查」，1740 万行的证明是「能过 Lean 的假推理」。
本报告\textbf{不}证明也不反证黎曼猜想；本报告只回答三个可以机械核查的问题：

\begin{enumerate}[label=\textbf{Q\arabic*.},leftmargin=3em]
\item \textbf{事实}：Nielsen 关于代码的各项陈述是否属实？特别是，这个恒真谓词是否真的位于主定理的依赖链上？
\item \textbf{逻辑}：如果属实，能否由此推出「该证明无效／空洞」？
\item \textbf{独立排查}：撇开 Nielsen 的论证，这份证明是否存在任何\emph{真正的} Lean reward hacking（额外公理、\lean{sorry}、偷换命题、绕过内核等）？
\end{enumerate}

\subsection{结论}

\begin{verdictbox}
\begin{enumerate}[label=\textbf{A\arabic*.},leftmargin=3em]
\item \textbf{事实基本属实，且比公开分析所确认的更强。}\lean{anyGeometry} 的定义逐字属实（24 个文件、118 行引用）。
我们对官方公布的 61.57\,GB 内核级证明导出文件做了逐声明的依赖传播，确认：五个主定理全部依赖 \lean{anyGeometry}
（依赖锥共 1,468,418 个声明，其中恰有 30 个依赖它），也全部依赖代码自称退化的几何模型 \lean{geoDegenerate} 等。
唯一的细节出入是：Nielsen 点名的 \lean{plugAll} 本身\emph{不在}依赖锥内，真正在路径上的是与之平行的孪生定理
\lean{every_character_of_green_v_win}（实质相同）。
\item \textbf{推理不成立。}「中间定义平凡 $\Rightarrow$ 证明无效」是一个不合逻辑的推断（non sequitur）：
  \begin{itemize}
  \item \emph{一般层面}：\lean{riemannHypothesis} 是无前提的闭定理，类型就是 Mathlib 的常量 \lean{RiemannHypothesis}。
  「空洞为真」只可能发生在\emph{带前提}的命题上；定义只是缩写，不可能凭空增加可证命题。据公开转述，Nielsen 的论文也承认该发布能对着 Mathlib 的 \lean{RiemannHypothesis} 通过类型检查。
  \item \emph{具体层面}：几何谓词 $\Gamma$ 在主路径上是「只写不读」的槽位——所有接收它的引理都对任意 $\Gamma$ 成立；
  封装它的归纳类型 \lean{GeoCl} 的消去子（recursor）在整个依赖锥中\emph{从未被任何声明引用}；最终假设包 \lean{HypsAlt} 的类型根本不含 $\Gamma$。
  把一个从不被读取的槽位填成 \lean{True}，等价于删除这个槽位，对结论没有任何影响。
  \end{itemize}
\item \textbf{未发现真正的 reward hacking。}源码中 \lean{axiom} 声明 0 处；\lean{sorry} 仅 5 处且全部在社区后加、不属于提交内容的 \path{Challenge.lean}；
\lean{native_decide}、\lean{implemented_by}、\lean{unsafe}、\lean{opaque}、\lean{extern} 均为 0；
导出文件中的公理恰为 \lean{propext}、\lean{Quot.sound}、\lean{Classical.choice} 三条；
没有对 \lean{RiemannHypothesis}、\lean{riemannZeta} 等名字的重定义或记号劫持；Mathlib 钉在官方标签 \texttt{v4.33.1}。
\conlecheSummary
\end{enumerate}
\end{verdictbox}

\subsection{逐条核查表}

Nielsen 论文正文受 PhilArchive 的反爬验证保护，本次审计未能取得全文（见第~\ref{sec:limits}~节）；下表所列主张取自其摘要被公开转述的句子、
其本人在 X 上的表述，以及 DEV Community 上对该文的逐项复核文章。

\begin{center}
\small
\renewcommand{\arraystretch}{1.25}
\begin{tabularx}{\textwidth}{@{}c>{\raggedright\arraybackslash}p{3.7cm}>{\centering\arraybackslash}p{1.35cm}>{\raggedright\arraybackslash}X>{\raggedright\arraybackslash}p{2.35cm}@{}}
\toprule
编号 & Nielsen 的主张 & 事实核查 & 依据（本报告章节） & 能否推出「证明无效」\\
\midrule
C1 & 几何谓词被定义为 \lean{True} & \OK & \path{OverIx.lean:19} 逐字属实（\S\ref{sec:fact-def}） & \NOTFOLLOW \\
C2 & 它被传入构造最终假设包的各处，并经 \lean{plugAll} 得到全体特征的结论 & \PARTLY & 确在依赖锥内（内核级确认）；但 \lean{plugAll} 不在锥内，实际路径是其孪生定理（\S\ref{sec:fact-path}） & \NOTFOLLOW \\
C3 & 几何见证由代码自称 degenerate 的模型构造 & \OK & \lean{geoDegenerate}、\lean{mdlJunk} 等确在锥内（\S\ref{sec:fact-degenerate}） & \NOTFOLLOW \\
C4 & 格见证「按定义」满足其性质 & \PARTLY & \texttt{nprimeDef := fun \_ => rfl} 属实；其余格性质是实证的（\S\ref{sec:fact-rfl}） & \NOTFOLLOW \\
C5 & 高度方程由 \lean{rfl} 成立，因为高度函数被赋值为自己的公式 & \OK & \texttt{h\_eq := fun \_ => rfl}（\S\ref{sec:fact-rfl}） & \NOTFOLLOW \\
C6 & 能从几何见证中提取内容的引理从未被调用 & \OK & 属实，且可加强为内核级结论：\lean{GeoCl.rec} 在锥中零引用（\S\ref{sec:fact-unused}） & \NOTFOLLOW；这恰恰是「死参数」的证据 \\
C7 & 因此该发布未被证明包含 RH 的任何证明 & \NO & 见第~\ref{sec:logic}~节 & —— \\
\bottomrule
\end{tabularx}
\end{center}

\subsection{怎样读这个结论}

\begin{itemize}
\item 本报告的结论是\textbf{条件于可信基}的：Lean 的类型论及其（多个相互独立的）内核实现是可靠的；Mathlib 对 $\zeta$ 与 \lean{RiemannHypothesis} 的定义表达了通常意义下的黎曼猜想；
我们核对过哈希的导出文件确由已公布的源码生成（最后一点我们依赖 MIT 方面的沙箱重建，未能在 7 小时内自行重建 17.5 万个模块，见第~\ref{sec:limits}~节）。
\item 本报告\textbf{没有}阅读并理解这 1740 万行证明的数学内容，也不对「这个证明的数学思想是什么」发表意见。
Nielsen 的观察确实揭示了一件值得解释的事：库中那套「真几何」谓词（\lean{geoOf}、\lean{IsGeoOf}、\lean{HasTrueSurf}）完全不在最终定理的依赖锥内。
这是一个关于\emph{如何理解与讲述}这份证明的合理问题，但它不是一个关于证明\emph{是否有效}的问题。
\end{itemize}

\section{背景、审计对象与来源完整性}

\subsection{事件时间线}

\begin{itemize}
\item \textbf{11 月 20 日}：Levent Alpöge 在 \url{https://alpo.ge/remark.pdf} 宣布与 Claude 合作得到黎曼猜想（以及 Dirichlet $L$-函数的广义黎曼猜想）的证明，并且\emph{只}公布了 Lean 形式化代码
（14 个分卷 \path{lean.tar.gz.part01}–\path{part14}），没有公布人类可读的论文。公告原话是 ``Claude and I have also formalized the arguments without assumptions''。
\item \textbf{11 月 21 日}：Lean 与 Mathlib 社区发表联合声明，认定陈述与证明均正确；MIT 的 Andrew Sutherland 公布了验证文件页，包括 61.57\,GB 的证明导出文件及多个独立内核的检查日志。
\item \textbf{11 月 21–22 日}：Jenny Lorraine Nielsen 在 PhilArchive 发布并多次更新《Invalidity of the Alpoge Proposed Proof of the Riemann Hypothesis》，
核心论据是几何谓词被定义为 \lean{True}，构造最终假设包的见证「其所需性质是空洞成立的」，结论是该发布「未被证明包含黎曼猜想的任何证明」。
\item \textbf{11 月 22 日}：DEV Community 上一篇复核文章确认了 \lean{anyGeometry} 定义的存在，但明确表示
\emph{没有}完成从 \path{Main.lean} 到 \lean{anyGeometry} 的依赖追踪，也没能找到「代码自称 degenerate 的模型」。这正是本报告补上的两处空白。
\end{itemize}

\subsection{审计对象与完整性核对}
\label{sec:provenance}

审计的可信度取决于「我们看的东西就是作者发布的东西」。下表中每一项哈希都是\emph{在本机重新计算}后与公布值比对的。

\begin{center}\small
\renewcommand{\arraystretch}{1.2}
\begin{tabularx}{\textwidth}{@{}>{\raggedright\arraybackslash}p{3.5cm}>{\raggedright\arraybackslash}X>{\raggedright\arraybackslash}p{4.3cm}@{}}
\toprule
对象 & 来源与 SHA-256（本机计算） & 核对结果 \\
\midrule
作者原始源码包 \path{lean.tar.gz}（239,275,554 字节） & 从 \path{alpo.ge} 下载 14 个分卷后拼接；\texttt{9631b4f0…93f49c7fc} & 与 MIT 页面公布的哈希一致 \\
源码树（175,187 个 \path{.lean} 文件） & 解包上述源码包；与 GitHub 镜像 \path{the-omega-institute/riemann-hypothesis}（提交 \texttt{c7d25b3b}）做 \texttt{diff -rq} & \path{src/} 逐文件一致；镜像仅多出社区添加的 \path{Challenge.lean}（及 \path{lakefile.toml} 末尾 5 行） \\
证明导出 \path{solution.ndjson.zst}（7,699,933,588 字节） & 从 MIT 验证页下载；\texttt{61a8938a…a681b05381} & 与公布值一致 \\
解压后的 \path{solution.ndjson}（61,570,595,788 字节） & \texttt{5ea74b91…f6be244e} & 与公布值一致 \\
挑战命题导出 \path{challenge.ndjson}（350,862,311 字节） & 从 MIT 验证页下载；\texttt{bd484ac9…bdd6ab2b9} & 与公布值一致；\textbf{且与我们在本机用官方 Mathlib 独立重建的文件逐字节相同}（\S\ref{sec:rebuild}） \\
Mathlib 版本 & \path{lakefile.toml} 钉住 \texttt{0df444a360ea…} & \texttt{git ls-remote} 证实它就是官方仓库的标签 \texttt{v4.33.1} \\
\bottomrule
\end{tabularx}
\end{center}

导出目标（\path{export-targets.txt}）是五个主定理加上内核原语，因此导出文件\emph{恰好}是这五个定理的传递依赖闭包（依赖锥）：
一个声明出现在导出文件里，当且仅当至少一个主定理（直接或间接）用到了它。这一点是第~\ref{sec:facts}~节内核级结论的基础。

\subsection{使用的工具}

\begin{itemize}
\item \textbf{源码检索}：\texttt{git grep}／\texttt{ripgrep} 全库检索，以及一份一次性生成的声明头索引（1,061,434 条）。
\item \textbf{conescan}（本次审计自写，约 320 行、无第三方依赖的 Rust 程序，源码随报告提交）：对 lean4export 的 NDJSON 流做单遍前向扫描，
按表达式 DAG 做污点传播，精确算出「哪些声明传递依赖某个种子常量」，同时给出公理清单、声明计数和见证链。
\item \textbf{con-leche}（\path{leanprover/con-leche}，提交 \texttt{bc667610}）：Lean FRO 的独立检查器，带有相对于集合论的一致性证明；我们在本机用官方 Lean \texttt{v4.36.0} 工具链从源码编译后运行。
\item \textbf{官方工具链}：Lean \texttt{v4.33.1}／\texttt{v4.36.0} 发行包、\path{lean4export}（标签 \texttt{v4.33.0}，提交 \texttt{15f6055e}）、Mathlib 官方缓存。
\item \textbf{交叉验证}：4 个 \texttt{agy} 代理（Gemini 4 Neon, XHigh）在只读条件下各自独立完成一项子审计（依赖追踪、$\Gamma$ 的逻辑角色、次要主张、卫生扫描），报告原文见仓库 \path{evidence/agy/}。
凡被本报告采用的代理结论，都由主审计方回到源码或内核级导出上复核过；仅由代理给出而未复核的内容会明确标注。
\end{itemize}

\section{Q1 事实核查：Nielsen 说的代码事实是否属实}
\label{sec:facts}

本节所有源码摘录均由脚本从\emph{作者原始源码包}（哈希已核对）按行号截取，未经手工改写；过长的行在报告中截断并标注。

\subsection{C1：恒真的几何谓词——属实}
\label{sec:fact-def}

\excerpt{anygeo}

同一个文件里并排放着两个同类型（\lean{GeoP}）的谓词：\lean{geoOf} 要求存在正整数 $c$ 使 \lean{IsGeoOf} 成立，并且 \lean{HasTrueSurf} 成立；
\lean{anyGeometry} 对任何输入都返回 \lean{True}。全库共有 24 个文件、118 行出现 \lean{anyGeometry}（与外部报道一致）。

谓词 $\Gamma$ 通过下面两个结构进入证明：\lean{OAt} 是带 \lean{geo}、\lean{lam} 两个字段的「证书」，\lean{GeoCl} 是把 $\Gamma$ 沿除子变换封闭起来的归纳谓词。

\excerpt{oat}
\excerpt{geocl}

\subsection{C2：它在主定理的依赖链上吗——在；但点名的 \texttt{plugAll} 不在}
\label{sec:fact-path}

这是公开讨论中悬而未决的关键问题（DEV 文章原话：没有确定 \path{Main.lean} 的依赖路径是否经过 \lean{anyGeometry}）。我们从两个独立的层面回答。

\paragraph{（a）源码层面。}主定理文件只有 12 行：

\excerpt{main}

沿着 \lean{rhOfAll} $\to$ \lean{allChi} $\to$ \lean{UpperTheory.FinalTheorems.everyCharacter} $\to$ \lean{…every_character_of_literature_of_vWin}
一路向下，到达下面这个定理；它的证明三次把 \texttt{geoCl anyGeometry} 作为显式参数传给下游引理：

\excerpt{siegel_a}
\excerpt{siegel_b}

所以 \lean{anyGeometry} 确实在源码层面的调用链上。但请注意 Nielsen 点名的定理 \lean{plugAll}：

\excerpt{plugall}
\excerpt{everyE2}

全库检索表明，\lean{plugAll} 只有一个调用者 \leanu{everyE₂}，而 \leanu{everyE₂} \emph{没有任何调用者}。主路径实际使用的是 \path{ThirdPlaceSeries/} 目录下与之逐句平行的孪生版本
（\lean{HypsAlt} 对应 \lean{Hyps}，\lean{every_character_of_green_v_win} 对应 \lean{plugAll}，\lean{parametersOfExtension} 对应 \lean{prmR}）：

\excerpt{topplug}

两条平行路线同样硬编码了 \texttt{geoCl anyGeometry}，因此这一出入不影响 C2 的实质，只影响其表述的精确性。

\paragraph{（b）内核层面（决定性）。}源码检索只能给出「文本上引用了」的近似结论。为得到精确答案，我们用 conescan 扫描了官方导出文件
（1,080,769,508 行、1,065,214,335 个表达式结点、1,468,418 条声明记录），以 \lean{GlobalTheory.Mem.Top.Bl.anyGeometry} 为种子做污点传播。结果：

\begin{verdictbox}
在五个主定理的依赖锥（1,468,418 个声明）中，\textbf{恰有 30 个声明}传递依赖 \lean{anyGeometry}，五个主定理全部在内。一条内核级见证链是：

{\small\lean{riemannHypothesis} $\to$ \lean{rhOfAll} $\to$ \lean{allChi} $\to$ \lean{UpperTheory.FinalTheorems.everyCharacter} $\to$ \lean{…everyCharacter_of_row}
$\to$ \lean{UpperTheory.LiteratureTheorems.every_character_of_literature_of_vWin} $\to$ \lean{UpperTheory.Lit.every_character_of_theta_integrals_local_matching}
$\to$ \lean{UpperTheory.Lit.nonempty_hyps_definite_side_of_closed_bound} $\to$ \lean{GlobalTheory.Fc.S0.StageU.hypotheses_nonempty_of_unramified_exponents_alt}
$\to$ \lean{GlobalTheory.Mem.Top.Bl.anyGeometry}。}

同一次扫描还表明：\lean{UpperTheory.Top.plugAll}、\lean{plugWin}、\leanu{everyE₂}、\leanu{hypsE₂}、\lean{GlobalTheory.Fc.Hyps}、\lean{prmR}、\lean{Prm}
这些名字\textbf{从未出现在导出文件中}，即它们不在任何主定理的依赖锥内。
\end{verdictbox}

30 个声明的完整清单见附录表~\ref{tab:tainted}。值得注意的是其中「出现在类型里」一列：从 \lean{every_character_of_theta_integrals_local_matching} 往上，
\lean{anyGeometry} 只出现在\emph{证明项}里而不出现在\emph{命题}里——这些定理的陈述已经与 $\Gamma$ 无关。

\subsection{C3：「代码自称 degenerate 的模型」——属实}
\label{sec:fact-degenerate}

DEV 文章称其检索 \texttt{degenerate} 只找到数学术语 nondegenerate，因而未能确认这一条。实际上相关定义采用驼峰命名（大写 D），区分大小写的检索会漏掉它：

\excerpt{geodegen}

这是抽象几何接口 \texttt{DV.Geo F} 的一个「占位」实现：簇取为 $\mathbb{N}$，有效性、nef、big 等谓词一律取 \lean{True}，闭子簇类取单点 \lean{PUnit}。
它所在的文件就叫 \path{Inhabited.lean}，同文件还有 \lean{cxJunk}、\lean{mdlJunk}、\lean{algDegenerate}、\lean{padDegenerate} 等同类定义。
它们经由 \lean{divXFl} 进入主路径上的引理 \lean{geoOver_stages}：

\excerpt{divxfl}
\excerpt{geoover}
\excerpt{geoover2}

第二遍 conescan 以这些名字为种子，确认它们\emph{全部}位于 \lean{riemannHypothesis} 的内核级依赖锥中：

\begin{center}\footnotesize
\begin{tabularx}{\textwidth}{@{}l>{\raggedright\arraybackslash}Xrc@{}}
\toprule
种子 & 导出文件中匹配到的全名 & 依赖它的声明数 & 主定理依赖它？\\
\midrule
\tableinput{tables/scan2_counts.tex}
\bottomrule
\end{tabularx}
\end{center}

{\small 注：\lean{geoDegenerateNoTau} 的名字在导出中出现但没有对应的声明记录、也没有被任何声明引用；\lean{geoTrue}（由「真几何」\lean{geoOf} 构造的 \lean{DV.Geo}）、
\lean{shimUJunk}、\lean{NormalizedHeckeDataPlaceholder} 的名字完全没有出现。末两行是与第一遍扫描的一致性对照。}

\paragraph{与 C1 的一个重要区别。}第三遍扫描表明，抽象几何接口 \lean{Formalism.BasicNotions.DV.Geo} 与 $\Gamma$ 槽的处境并不相同：
依赖锥中有 736 个声明（占 1,468,418 的 0.05\%）以某种方式涉及 \lean{DV.Geo} 类型，其中 621 个（传递地）读取了它的某个字段——
实际被用到的字段投影只有 \lean{Var}、\lean{Dv}、\lean{H0}、\lean{inter}、\lean{vanishOn}、\lean{arVol} 等 18 个。
也就是说，这个接口\emph{不是}无人读取的死槽，而是「一小层对任意几何数据成立的泛型声明，被用在了一个平凡实例上」。
这不改变第~\ref{sec:logic}~节的结论，但两种情形应当分开表述，本报告不把「只写不读」这个更强的说法用在 \lean{geoDegenerate} 上。

\subsection{C4、C5：「按定义成立」与「由 \texttt{rfl} 成立」——机制属实}
\label{sec:fact-rfl}

\excerpt{heighteq}

这段代码先把高度函数\emph{定义}为显式公式 \texttt{ht := DV.hExpl6 …}，再用 \lean{setDiv} 构造一个以它为高度的新数据，于是「高度等于该公式」（\lean{h_eq}）和
「格等于 \lean{nPrime}」（\lean{nprimeDef}）自然由 \lean{rfl} 成立。\lean{LedPart.pt}、\lean{LedPart.ht}、\lean{setDiv} 均在依赖锥内（第二遍扫描确认）。
所以 C5 的描述在机制上是准确的；C4 则只对 \lean{nprimeDef} 这类字段准确——交叉验证代理 C 指出同一见证的其他格性质（如 \lean{NuDef}、\lean{nuLat_comap_isLattice}）是用实质性定理证明的，
这一点我们只做了抽查，未逐项复核。

「先定义对象、再由定义读出它满足的等式」是构造性存在证明的标准写法：要证 $\exists h,\ h = F$，取 $h := F$ 即可。它是否「掏空」了什么，取决于下游是否需要 $h$ 是某个\emph{独立定义}的量——
而这正是第~\ref{sec:logic}~节要回答的问题。

\subsection{C6：「提取几何内容的引理从未被调用」——属实，且可加强为内核级结论}
\label{sec:fact-unused}

源码检索显示：\lean{GeoCl} 这个名字在全库只出现在它的定义文件 \path{NotGeoClFalse.lean} 和一处 \lean{GeoCl.step} 的使用中；
构造子 \lean{GeoCl.base} 只经由引理 \lean{geoCl_of} 被用到一次；全库没有任何写明针对 \lean{GeoCl} 的分情形分析或归纳。
「真几何」谓词 \lean{GlobalTheory.Mem.Top.Bl.geoOf} 只在定义处和 \path{NotLamAreShift.lean} 的三个无人调用的引理中出现
（按单词匹配 \lean{geoOf} 共 10 个文件，其余是不同命名空间下的同名定义）。

但「文本上没有」不等于「证明项里没有」（模式匹配可以不写出类型名）。内核级扫描给出了无歧义的答案：

\begin{verdictbox}
\begin{itemize}
\item \textbf{\lean{GeoCl.rec} 零引用。}在归纳命题上做任何分情形分析或归纳，展开到内核层面都必须经过它的递归子。
导出文件中引用 \lean{GlobalTheory.Mem.Top.Bl.GeoCl.rec} 的只有 \lean{GeoCl} 自己的归纳声明块（4 条记录），辅助定义 \lean{GeoCl.casesOn}、\lean{GeoCl.recOn} 的名字根本没有出现。
也就是说，\emph{依赖锥里没有任何一个证明曾经「打开」过一个 \lean{GeoCl} 证书}。
\item \textbf{\lean{OAt} 的 \lean{geo} 字段只被两个定理读取}（\lean{divEx_geoCl} 与 \lean{oAt_setDiv}），而这两处读取的目的都只是把它原样塞进 \lean{GeoCl} 的构造子里（见下）。
结构 \lean{OAt} 的递归子 \lean{OAt.rec} 同样零引用，所以投影是读取其字段的唯一途径。
作为对照，同一结构的另一个字段 \lean{lam} 除了这两处转手之外，还被 \lean{GlobalTheory.T.parameter_of_discrete_series_ray_alt} 真正使用——
同一个证书里，一个字段有消费者，另一个没有。
\item \textbf{\lean{geoOf}、\lean{IsGeoOf}、\lean{HasTrueSurf} 不在依赖锥内}——这三个名字在导出文件中完全没有出现。
\end{itemize}
\end{verdictbox}

\excerpt{divexprov}
\excerpt{divexgeocl}
\excerpt{oatsetdiv}

\lean{divEx_prov} 是整条路径上真正「兑现」$\Gamma = \mtt{True}$ 的地方：它用 \lean{trivial} 填上 \lean{geo} 字段，同时用实证的 \texttt{lamAre\_heq …} 填上 \lean{lam} 字段。
之后 \lean{geo} 只被转手（\lean{GeoCl.base}、\lean{GeoCl.step}），再也没有被打开。

\subsection{小结}

Nielsen 的事实陈述经得起核查，有的地方（C3、C6）甚至比已有的公开复核所能确认的更扎实。争议的全部重量因此落在下一个问题上：这些事实\emph{意味着}什么。

\section{Q2 逻辑核查：由这些事实能否推出「证明无效」}
\label{sec:logic}

\subsection{被审计的推理}

把 Nielsen 的论证整理成标准形式：

\begin{notebox}
\begin{enumerate}[label=(P\arabic*),leftmargin=3em]
\item 几何谓词被实例化为恒真谓词，几何见证来自一个退化模型，若干等式由 \lean{rfl} 成立。\hfill（事实，属实）
\item 这些见证进入了构造最终假设包的过程。\hfill（事实，属实）
\item 所以见证「所需的性质是空洞成立的」，见证「是废的」。\hfill（解释）
\item 所以该发布并不包含黎曼猜想的证明。\hfill（结论）
\end{enumerate}
\end{notebox}

从 (P1)(P2) 到 (P3)(P4) 这一步是全部论证的承重点。它不成立，原因有一般性的，也有针对这份代码的具体原因。

\subsection{一般性原因：闭定理没有「空洞为真」的余地}

\paragraph{命题 1（定义是保守的）。}在 Lean 的类型论里，\texttt{def c : T := v} 只是给一个已经通过类型检查的闭项 $v$ 起名字。
任何用到 $c$ 的证明项，把 $c$ 全部展开（$\delta$-约化）之后仍然是类型不变的合法证明项。因此「某个中间定义取了平凡的值」\emph{不可能}让原本不可证的命题变得可证。
把证明里所有的 \lean{anyGeometry} 换成 \texttt{fun \_ \_ \_ => True}，得到的是同一个证明。

\paragraph{命题 2（空洞为真只属于带前提的命题）。}「空洞成立」说的是形如 $H \to P$ 的命题：当 $H$ 不可满足，或 $H$ 被人为弱化到与 $P$ 无关时，证明 $H\to P$ 并不能说明 $P$。
而本案的最终定理是
\begin{Verbatim}[fontsize=\small,frame=single,rulecolor=\color{gray!50}]
theorem riemannHypothesis : RiemannHypothesis := rhOfAll
\end{Verbatim}
\noindent 它没有任何前提；内核级扫描确认其类型就是常量 \lean{RiemannHypothesis}，而 con-leche 的比对（\S\ref{sec:conleche}）表明这个常量及其读到的全部声明，与只导入 Mathlib 的挑战导出逐条相同。
中间引理当然有前提，但每一个前提在被使用的那一刻，都必须由一个内核检查过的证明项来满足——「前提被满足得很轻松」不是「前提没有被满足」。

\paragraph{命题 3（内核接受意味着什么）。}如果内核接受了 $t : \mtt{RiemannHypothesis}$，且 $t$ 的依赖只含三条标准公理，那么要否定「RH（按 Mathlib 的表述）已被证明」，逻辑上只剩四条路：

\begin{enumerate}[label=(\roman*)]
\item 证明用了额外公理或 \lean{sorry}；
\item 被证明的命题不是黎曼猜想（名字被劫持、Mathlib 被改动、定义本身有误）；
\item 内核实现有错误，而且被这份证明利用了；
\item 被检查的产物不是声称的产物（导出文件与源码不对应）。
\end{enumerate}

Nielsen 的论证不属于其中任何一条，而且据公开转述，其论文自己也承认该发布能对着 Mathlib 的 \lean{RiemannHypothesis} 通过类型检查。
一旦承认这一点，「见证是否平凡」在逻辑上就与「定理是否成立」无关了。(i)–(iv) 四条路我们在第~\ref{sec:hygiene}、\ref{sec:replication}~节逐一排查。

\subsection{为什么直觉会出错：把谓词改弱，对生产者和消费者的作用相反}

「把该证的条件改成 True 就能蒙混过关」这个直觉，在只有\emph{一方}被检查时是对的。但一个谓词在证明里总是同时面对两方：

\begin{itemize}
\item \textbf{生产者}——要证明谓词成立的一方。谓词越弱，它越轻松；弱到 \lean{True}，只需写 \lean{trivial}。
\item \textbf{消费者}——要\emph{利用}谓词成立来推出别的东西的一方。谓词越弱，它越艰难；弱到 \lean{True}，它从中什么也得不到。
\end{itemize}

作弊要得逞，必须让生产者占到便宜而消费者不付出代价。在 Lean 里这只有三种做法：让消费者不存在（把弱化的谓词直接留在\emph{最终}定理的陈述里），
让消费者不受检查（\lean{axiom}／\lean{sorry}），或者换掉最终的陈述。只要最终陈述固定、且没有公理，消费者就必须在\emph{没有}该谓词内容的条件下把结论证出来——
那么这个谓词本来就是多余的。

随报告附带的 \path{evidence/toy/ToyGeometry.lean} 用不到一百行代码复现了这个结构，并在 Lean 4.36.0 下编译通过：
\begin{itemize}
\item \textbf{Part A}：一个带 \lean{geo} 槽位的证书、一个对任意 $\Gamma$ 成立的消费者引理、一个不含 $\Gamma$ 的假设包、一个用 \lean{trivial} 填槽位的生产者。
最终的闭定理成立，\texttt{\#print axioms} 报告它\emph{不依赖任何公理}。
\item \textbf{Part B}：一旦消费者真的需要谓词的内容（需要 $\forall n,\ \Gamma\,n \to 0<n$），取 $\Gamma := \mtt{True}$ 会让这个条件直接变成假命题——我们给出了它的否定的证明。没有免费的午餐。
\item \textbf{Part C}：真正的作弊长什么样——额外公理（被 \texttt{\#print axioms} 揭穿）、把前提留在最终定理里、偷换目标命题。
\end{itemize}

\subsection{具体原因：在这份代码里，几何槽是「只写不读」的死参数}

一般性论证已经足够，但我们仍然逐项核实了 $\Gamma$ 在这份代码里的实际角色，因为它直接回答「那它为什么在那里」。

\paragraph{证据 E1：$\Gamma$ 是参数记录里的一个字段，消费它的引理全部对任意 $\Gamma$ 成立。}

\excerpt{params}
\excerpt{prm1}

主路径上接收 \texttt{geoCl anyGeometry} 的引理都把 $\Gamma$ 当作\emph{任意}参数。例如单调性引理与「几何弱相容」引理的陈述（节选）：

\excerpt{mono_sig}
\excerpt{geo0_sig}

内核级扫描也印证了这一点。下列声明在导出文件中\emph{不}依赖 \lean{anyGeometry}，也就是说它们是在不知道 $\Gamma$ 是什么的情况下被证明或定义的：
\begin{itemize}
\item \lean{GlobalTheory.Glob.open_and_linked_of_stronger_extension_alt}；
\item \lean{GlobalTheory.T.linkedAt_geometricWeakTies_of_tie}；
\item \lean{GlobalTheory.Mem.Top.Bl.parametersOfExtension} 与 \lean{GlobalTheory.Mem.Top.Bl.OX}。
\end{itemize}

\paragraph{证据 E2：下游真正需要的那条「几何」性质是被证明出来的，不是被假设出来的。}
主路径需要的十二条「开放参数陈述」里有一条名为 \lean{geo}：

\excerpt{openstmt}

它在 \lean{anyGeometry} 参数下的证明并没有用到 $\Gamma$ 的内容，而是由除子条件（\lean{hDiv}）、联结条件（\lean{hTie}）和射线类域论推出：

\excerpt{openorbit}

\paragraph{证据 E3：最终假设包的类型里没有 $\Gamma$。}
最终被 \lean{every_character_of_green_v_win} 消费的假设包是 \lean{HypsAlt}，一个含四十余个字段的结构（节选）：

\excerpt{hypsalt}

它的类型不含 \lean{GeoP}。内核级扫描确认 \lean{GlobalTheory.Fc.HypsAlt}、\lean{every_character_of_green_v_win}、\lean{zero_free_all_of_green_v_win} 均不依赖 \lean{anyGeometry} 或 \lean{GeoCl}：
从假设包到「全体 Dirichlet 特征在 $\operatorname{Re} s>1/2$ 无零点」的那一段推导，是在一个完全看不见几何槽的世界里完成的。

\paragraph{证据 E4：槽里的东西从未被读出。}
如 \S\ref{sec:fact-unused} 所述，\lean{GeoCl.rec} 在依赖锥中零引用；\lean{OAt} 的 \lean{geo} 字段的两处读取都只是把它转交给 \lean{GeoCl} 的构造子。
对任意 $\Gamma$ 证明的引理按参数性无法窥视 $\Gamma$ 的内容；取 $\Gamma=\mtt{True}$ 时本就无内容可读；取 $\Gamma = \mtt{GeoCl}\,(\cdots)$ 时又从未被消去。三种情形合起来就是全部情形。

\begin{verdictbox}
\textbf{结论。}在通往 \lean{riemannHypothesis} 的整条路径上，几何谓词 $\Gamma$ 是一个只被写入、被搬运、从未被读取的槽位。
把它实例化为 \lean{True} 与把这个字段从结构里删掉是等价的；对一个不被读取的槽位而言，\lean{True} 恰恰是最诚实的填法
（换成「真几何」\lean{geoOf}，反而要额外证明一批下游根本不用的几何事实）。
这是代码里的一个\emph{死参数}，不是证明里的一个\emph{漏洞}。

退化模型 \lean{geoDegenerate} 和「由 \lean{rfl} 成立」的等式（C3–C5）情形略有不同，结论相同。那里的抽象几何接口 \lean{DV.Geo} 并非无人读取（\S\ref{sec:fact-degenerate}）：
一小层泛型声明是对\emph{任意}满足接口约束的几何数据证明的，证明库随后挑了一个平凡的实例代入。这是「给抽象接口选一个具体模型」的标准做法；
\emph{这个模型是否足以推出关于真实 $L$-函数的结论}，由内核逐步检查，不取决于模型的名字好不好听。
如果一个平凡模型不足以支撑结论，后果是下游某个引理证不出来、内核报错，而不是一个通过检查的假定理。
\end{verdictbox}

\subsection{公平地说：Nielsen 的观察确实说明了什么}

\begin{itemize}
\item 它说明这份形式化证明\textbf{没有}通过这个接口使用任何非平凡的几何输入：\lean{geoOf}、\lean{IsGeoOf}、\lean{HasTrueSurf} 以及围绕它们的一批引理写在库里，却不在最终定理的依赖锥内。
最自然的解读是它们属于某条早期或备选路线留下的脚手架；全库有约三千个以 \texttt{\_alt}／\texttt{\_alt\_alt} 结尾的平行版本声明，也指向同一个方向。
\item 它提醒读者：\emph{名字不是证据}。看到 \path{…/AbelianVarietyTheory/…}、\lean{Geo}、\lean{HasTrueSurf} 这样的名字，并不能推断证明「用了」相应的数学；
同理，看到 \lean{anyGeometry}、\lean{geoDegenerate}、\lean{mdlJunk} 这样的名字，也不能推断证明「是假的」。能下结论的只有项和类型。
\item 它指向一个真实而尚未解决的问题：作者只发布了 Lean 代码，并自述「在完全掌握其周边数学之前就拿出了完整论证」。这份证明的数学思想究竟是什么、哪些部分是本质的，
目前没有人类可读的说明。\textbf{这是一个关于理解的问题，不是一个关于有效性的问题。}
\end{itemize}

\subsection{关于「让模型造一条能过 Lean 的假推理链」的说法}

委托人转来的评论称：让任意模型「不用真的证明，只需要制造足够长的能过 Lean 的逻辑链假装在推理」，就能得到一个「证明黎曼猜想」的结果。
这个说法混淆了两件事。一条\emph{最终命题正确、不含额外公理}而又通过内核的推理链，按定义就是一个证明——内核不关心链条长短，也不存在「假装通过」。
模型能够造假的地方只有命题 3 中列出的 (i)(ii)(iv)：加公理、留 \lean{sorry}、换命题、换产物。这些全部可以机械地检查，而且本案中全部检查过了。

\section{Q3 独立排查：有没有真正的 reward hacking}
\label{sec:hygiene}

这一节不依赖 Nielsen 的论证，而是按命题 3 的四条路，逐一寻找「能让 Lean 接受一个错误或被偷换的定理」的手法。

\subsection{(i) 公理、\texttt{sorry} 与信任编译器的后门}

对全部 175,187 个源文件做单遍正则扫描（原始命中 173,881 条，已随报告提交统计）：

\begin{center}\small
\renewcommand{\arraystretch}{1.15}
\begin{tabularx}{\textwidth}{@{}>{\raggedright\arraybackslash}p{5.2cm}r>{\raggedright\arraybackslash}X@{}}
\toprule
模式 & 命中数 & 说明 \\
\midrule
\lean{axiom} 声明 & \textbf{0} & 全库没有任何公理声明 \\
\lean{sorry} & 5 & 全部在 \path{src/Challenge.lean} 第 19–27 行；该文件是社区为 Comparator 添加的，不在作者源码包内，也不被 \path{Main.lean} 导入 \\
\lean{sorryAx}、\lean{admit} & 0 & \\
\lean{native_decide}、\lean{ofReduceBool}、\lean{ofReduceNat}、\lean{trustCompiler} & 0 & 没有把编译器纳入可信基 \\
\lean{implemented_by}、\lean{extern}、\lean{unsafe}、\lean{unsafeCast}、\lean{lcProof}、\lean{opaque}、\lean{csimp} & 0 & （\lean{allowUnsafeReducibility} 是一个 \lean{set_option} 名，见下） \\
\texttt{decide +kernel} 等 & 151 & 由内核求值，不扩大可信基 \\
\bottomrule
\end{tabularx}
\end{center}

内核级导出给出同样的结论，而且不依赖对源码的文本匹配：conescan 统计导出文件中的声明记录为
\lean{thm} 1,246,110、\lean{def} 214,241、\lean{inductive} 7,969、\lean{opaque} 91、\lean{quot} 4、\textbf{\lean{axiom} 3}，
三条公理的全名恰为 \lean{propext}、\lean{Quot.sound}、\lean{Classical.choice}；标记为 \lean{unsafe}／\lean{partial} 的记录为 0；
\lean{sorryAx}、\lean{Lean.ofReduceBool}、\lean{Lean.ofReduceNat}、\lean{Lean.trustCompiler} 这些名字在导出中完全没有出现。
声明总数 1,468,418 与 MIT 页面上 con-leche 报告的数字一致。

\subsection{(ii) 命题有没有被偷换}

\begin{itemize}
\item \path{Main.lean} 的五个定理没有任何前提（文件见 \S\ref{sec:fact-path}）；项目级选项 \texttt{autoImplicit = false}，不存在隐式变量混入。
\item 全库没有名为 \lean{RiemannHypothesis}、\lean{riemannZeta}、\lean{DirichletCharacter}、\lean{completedRiemannZeta} 的声明（声明头索引检索为空）；
出现 \lean{RiemannHypothesis} 这个名字的 26 个文件（含社区添加的 \path{Challenge.lean}）都是在\emph{使用}它。有两处名为 \lean{LFunction} 的定义位于各自的命名空间内（\path{GJLocalZeta.lean}、\path{HeckeLFunction.lean}），
与 Mathlib 的 \lean{DirichletCharacter.LFunction} 不冲突。
\item 记号与宏：全库 1,695 处 \lean{notation} 全部是 \lean{local}，另有 1 处 \lean{scoped}；\emph{没有}全局记号。全局 \lean{macro} 共 22 处，全部是策略宏；
全局 \lean{elab} 共 4 处（见下）。因此不存在通过记号劫持改变 \path{Main.lean} 或 \path{Statements.lean} 解析结果的可能。
\item \path{Statements.lean} 只 \texttt{import Mathlib}；\path{lakefile.toml} 与 \path{lake-manifest.json} 把 Mathlib 钉在官方仓库的 \texttt{0df444a3…}，
我们用 \texttt{git ls-remote} 证实这正是官方标签 \texttt{v4.33.1}；没有路径依赖或覆盖，额外的编译参数只有 \lean{--tstack=400000}。
\item 最强的一项证据见 \S\ref{sec:rebuild}：我们在本机用官方工具链和官方 Mathlib 缓存独立生成了「只导入 Mathlib、陈述五个命题」的挑战导出，它与社区公布的文件逐字节相同；
con-leche 的 \lean{--compare} 模式再把解答导出与它逐条比对。
\end{itemize}

\subsection{(iii) 有没有绕过或利用内核}

值得逐一查看的元编程只有下面几处，它们才是「真正的作弊」可能藏身的地方：

\paragraph{\texttt{BuildAid/D1..D5.lean}：\texttt{run\_cmd Lean.modifyEnv}。}乍看可疑，实际只是往库建议（library suggestions）的名字黑名单里加命名空间，与内核无关：
\excerpt{buildaid}

\paragraph{\texttt{cpt\_copy\_thms}：批量复制声明。}它用的是标准的 \lean{Lean.addDecl}，每条被复制的声明都要重新经过内核检查：
\excerpt{cptcopy}

\paragraph{\texttt{kexactB} 与 \texttt{kcopy}：把类型检查完全交给内核。}这两个自定义策略（\lean{kexactB}：50 个文件、234 处，另经简写宏 \lean{kc%} 在 32 个文件出现 182 处；\lean{kcopy}：含定义在内 4 个文件、5 处）直接把一个项赋给目标元变量，
跳过精化器（elaborator）的类型核对；\lean{kcopy} 还会手工拼装原始投影项 \lean{Expr.proj}：
\excerpt{kcopy}

这是合法的工程手段（省掉精化器昂贵的定义等价检查），但它意味着这些证明的正确性\emph{只}由内核把关——如果内核在投影的类型检查上有漏洞，这里正是能被利用的位置。
社区的验证说明也指出，Lean \texttt{v4.33.1} 自带内核单独通过只能算有限证据，裁决应由更新的内核和独立检查器作出。
所以对这一条，源码阅读不能给出最终答案，必须靠独立内核重放——这正是 \S\ref{sec:conleche} 的内容。

\paragraph{选项。}全库 \lean{set_option} 涉及 33 种选项，绝大多数是 \lean{maxHeartbeats} 等资源限制与 linter 开关；没有任何 \lean{debug.*} 选项（如 \lean{debug.skipKernelTC}）。
\lean{allowUnsafeReducibility}（37 处）只影响精化器一侧的可约性标注，不影响内核。

\subsection{(iv) 被检查的产物是不是声称的产物}

见 \S\ref{sec:provenance} 的完整性核对：源码包、镜像、导出文件的哈希全部在本机复算并与公布值一致。
唯一没有由我们亲自完成的一环是「由源码重新编译 17.5 万个模块并导出」；这一环我们依赖 MIT 页面所述的沙箱重建（其导出与公布文件逐字节相同），并在第~\ref{sec:limits}~节如实列为局限。

\subsection{小结}

按命题 3 的四条路逐一排查，没有发现任何一条成立的迹象。交叉验证代理 D 独立完成的卫生扫描得出了相同结论，并补充了两类本节未展开的良性机制（32 处内联 \lean{run_tac}、1 处 \lean{@[positivity]} 扩展）。

\section{本机独立复现}
\label{sec:replication}

除了阅读源码，本次审计在一台 Apple Silicon Mac（16 核、128\,GB 内存）上独立完成了下面四项可复现的实验。全部脚本、日志与输出见仓库 \path{evidence/}。

\subsection{conescan：内核级依赖锥分析}
\label{sec:conescan}

\paragraph{方法。}lean4export 的 NDJSON 导出按依赖顺序给出名字、表达式结点（每个结点一行，子结点以编号引用）和声明记录。
conescan 单遍读取整个文件，为每个表达式结点维护一个 16 位污点掩码：引用了种子常量或已被污染的声明的 \lean{const} 结点被置位，复合结点取其子结点的并；
每条声明记录的类型与值（归纳类型则包括所有构造子类型、递归子类型与计算规则右端）决定该声明的掩码。
由于声明按依赖顺序出现，一遍扫描即可得到「传递依赖某种子」的\emph{精确}声明集合，并为每个被污染的声明记下一个直接见证，从而还原出依赖链。
种子可以是常量名，也可以是「对结构 $S$ 的第 $i$ 个字段的投影」这样的表达式模式。

\paragraph{自检。}(1) 统计到的行数、字节数与声明总数（1,468,418）和官方数字一致；(2) 以 \lean{every_character_of_green_v_win} 为种子时，传播正确地到达了全部五个主定理（这是源码层面已知的依赖）；
(3) 两遍独立扫描中 \lean{anyGeometry} 与 \lean{OAt.geo} 投影的计数相同（各 30）。

\paragraph{注意。}conescan 是依赖分析器，不是类型检查器；它回答「谁用了谁」，不回答「证明对不对」。后者由下一小节负责。

\paragraph{主要结果。}见 \S\ref{sec:fact-path}、\S\ref{sec:fact-degenerate}、\S\ref{sec:fact-unused} 以及附录~\ref{app:conescan}。\scanThreeSummary

\subsection{con-leche：对 61.57\,GB 导出证明的独立内核重放与命题比对}
\label{sec:conleche}

con-leche 是 Lean FRO 发布的独立检查器，自带「不接受 \lean{False} 的证明」的形式化一致性证明，使用自己的项表示，不依赖官方内核的 C++ 实现。
我们用官方 Lean \texttt{v4.36.0} 发行包从源码编译（提交 \texttt{bc667610}，593 个构建任务），然后运行
\begin{Verbatim}[fontsize=\small,frame=single,rulecolor=\color{gray!50}]
con-leche --verified --compare challenge.ndjson solution.ndjson
\end{Verbatim}
\noindent 其中 \lean{--verified} 是一致性定理所覆盖的模式；\lean{--compare} 先检查解答导出的全部声明，再检查挑战导出，
并要求挑战中每个以 \lean{sorry} 占位的定理在解答中有同名、同类型的常量，且「该类型读到的一切声明在两边逐条相同」。

\conlecheDetail

\subsection{独立重建挑战命题导出}
\label{sec:rebuild}

\lean{--compare} 的意义取决于挑战文件本身是否可信：它必须真的只是「官方 Mathlib + 五个命题陈述」。我们没有把这一点当作既定事实，而是自己重建了一遍：

\begin{enumerate}
\item 下载官方 Lean \texttt{v4.33.1} 发行包；
\item 建立一个只含 \path{Challenge.lean}（5 个带 \lean{sorry} 的定理，仅 \texttt{import Mathlib}）的工程，Mathlib 钉在 \texttt{0df444a3}；
\item 用 \texttt{lake exe cache get} 从 Mathlib 官方缓存取得 8,690 个预编译文件（不在本机编译 Mathlib）；
\item 用 \path{lean4export}（标签 \texttt{v4.33.0}）按 \path{export-targets.txt} 的目标导出。
\end{enumerate}

\begin{verdictbox}
重建得到的 \path{my_challenge.ndjson} 为 350,862,311 字节，SHA-256 为 \texttt{bd484ac9…bdd6ab2b9}，与社区公布的 \path{challenge.ndjson} \textbf{逐字节相同}（\texttt{cmp} 无差异）。
因此，用来比对的「标准答案」确实只由官方 Mathlib 和五个命题陈述决定，与证明库无关。
\end{verdictbox}

{\small 顺带一提：MIT 页面写该文件解压后为 350,862,970 字节，与我们得到的 350,862,311 字节相差 659 字节，但 SHA-256 与页面公布值一致，且与本机重建结果一致，应为页面上的笔误。}

\subsection{玩具模型}

\path{evidence/toy/ToyGeometry.lean}（95 行）在 Lean 4.36.0 下编译通过，输出为：
\VerbatimInput[fontsize=\small,frame=single,rulecolor=\color{gray!50}]{../evidence/toy/ToyGeometry.out}
\noindent 即：经由恒真谓词得到的闭定理 \lean{Toy.target} 不依赖任何公理；而真正作弊的 \lean{Toy.bad₁} 立刻被 \texttt{\#print axioms} 指认。

\subsection{交叉验证：四个独立的 \texttt{agy} 代理}

按委托人的建议，我们启动了四个 \texttt{agy --model gemini-4-neon-xhigh} 代理，在只读条件下各自独立审计一个子问题。
它们启动时主审计方尚未形成结论，提示词要求「对证明和对批评者都保持怀疑」；但需如实说明，代理 B 的提示词里带有「无前提的闭定理不可能空洞」这一提示，
因此它在这一点上的一致不能算完全独立。提示词原文与报告全文见 \path{evidence/agy/}。结果如下：

\begin{center}\small
\renewcommand{\arraystretch}{1.2}
\begin{tabularx}{\textwidth}{@{}c>{\raggedright\arraybackslash}p{3.3cm}>{\raggedright\arraybackslash}X@{}}
\toprule
代理 & 子问题 & 独立得出的结论（与主审计方的关系） \\
\midrule
A & 声明级依赖追踪 & \lean{anyGeometry} 在 \lean{riemannHypothesis} 的依赖锥内；\leanu{everyE₂}／\lean{plugAll}／\lean{Hyps} 不在；\lean{geoOf} 完全不在。\textbf{与 conescan 的内核级结果一致。} \\
B & $\Gamma$ 的逻辑角色 & 全部消费引理对任意 $\Gamma$ 成立；\lean{OAt.geo} 与 \lean{GeoCl} 没有提取内容的消去；最终假设包不含 $\Gamma$；「$\Gamma:=\mtt{True}$ 揭示的是一个残留的未用槽位而非逻辑缺陷」。\textbf{与第~\ref{sec:logic}~节一致。} \\
C & 次要主张 (iii)–(vi) & 首先指出 \lean{geoDegenerate}／\lean{mdlJunk} 等定义的存在（主审计方随后在源码与内核级导出上复核属实）；(v) 机制属实；(vi) 属实；(iv) 作为笼统表述不成立。 \\
D & 卫生扫描 & 0 \lean{axiom}、0 \lean{native_decide} 等；\lean{sorry} 仅在 \path{Challenge.lean}；无名字劫持；并指出 \lean{kexactB}／\lean{kcopy} 的存在（主审计方复核并写入 \S\ref{sec:hygiene}）。 \\
\bottomrule
\end{tabularx}
\end{center}

四个代理的结论彼此一致，也与主审计方一致：事实属实，推理不成立。需要说明的是，代理与主审计方读的是同一份源码，因此这是对\emph{阅读与推理}的交叉验证，而不是对源码真实性的独立来源。

\section{局限、结论与建议}
\label{sec:limits}

\subsection{本次审计没有做到的事}

\begin{enumerate}
\item \textbf{没有取得 Nielsen 论文的全文。}PhilArchive／PhilPapers 对自动访问返回反爬验证页面（HTTP 403），我们没有绕过它；查询时互联网档案馆也没有存档。
本报告核查的主张来自：该文摘要被多处转述的句子、作者本人在 X 上的表述、以及 DEV Community 复核文章所引用的内容。
如果论文正文（据报道已更新十余版）包含未被这些来源反映的其他论据，本报告没有评估它们。本报告的结论针对的是公开流传的那条论证——即委托人截图中的那条。
\item \textbf{没有从源码重新编译整个证明库。}17.5 万个模块的完整构建在单机上无法在 7 小时预算内完成（公开的验证使用了云端多机构建）。
「导出文件确由这份源码生成」这一环我们依赖 MIT 页面所述的沙箱重建（重建所得导出与公布文件逐字节相同）；我们自己核对的是源码包与导出文件各自的哈希，以及挑战导出的独立重建。
\item \textbf{没有审读数学内容。}我们没有、也不可能在此时间内读懂 1740 万行证明的数学。本报告对「证明的数学思想是否自然、是否可以简化、几何在其中扮演什么角色」不作判断。
\item \textbf{可信基。}结论依赖于：Lean 的类型论是一致的；多个独立内核实现（官方内核、nanoda、lean4lean、con-leche、con-ron）不会同时在同一处出错；
\path{lean4export} 忠实地导出了环境；Mathlib 对 \lean{riemannZeta} 与 \lean{RiemannHypothesis} 的定义表达了通常意义的黎曼猜想。
其中最后一项是数学共同体的判断，不是本报告能机械验证的；MIT 页面提到的 Robin 不等式与 Lagarias 不等式推论提供了一个不依赖 $\zeta$ 定义的旁证，我们未复核。
\item \textbf{工具本身。}conescan 是为本次审计临时编写的（源码约 320 行，已提交供复核），其正确性依赖于「导出按依赖顺序排列」这一假设和若干自检；
con-leche 是一个很新的工具，其实现由 AI 编写并附形式化证明。我们把二者与第三方已公布的 nanoda、lean4lean 结果互相印证，而不单独依赖其中任何一个。
\item \textbf{交叉验证的独立性有限。}四个 \texttt{agy} 代理与主审计方阅读的是同一份源码；被本报告采用的代理结论均经复核，未复核者已在文中标明。
\end{enumerate}

\subsection{结论}

\begin{verdictbox}
\textbf{Nielsen 的打假不成立。}

\begin{itemize}
\item 其\textbf{事实观察}大体准确，本报告还把其中最关键的一条（\lean{anyGeometry} 与退化模型确在主定理依赖链上）从「未确认」提升为「内核级确认」，并纠正了一处细节（\lean{plugAll} 不在链上，在链上的是其孪生定理）。
\item 其\textbf{推理}不成立：对一个不带前提、类型为 Mathlib \lean{RiemannHypothesis}、只依赖三条标准公理、并经多个独立内核检查的闭定理而言，
「中间见证平凡」既不构成也不蕴含「证明无效」。在这份代码里，被指为「造假」的几何槽是一个只写不读的死参数。
\item 我们按四条可能的造假途径逐一排查，\textbf{没有找到这份证明的任何 bug 或 reward hacking}。就委托人「找出证明的 bug 就算赢」的目标而言，本次审计的如实答案是：没有找到，而且有充分理由相信在被审计的范围内不存在。
\end{itemize}
\end{verdictbox}

\subsection{建议}

\begin{itemize}
\item \textbf{给读者：评估「Lean 证明造假」类说法的五问清单。}
(1) 最终定理的陈述是不是你以为的那个命题（最好由只导入标准库的独立文件陈述并机械比对）？
(2) 它有没有前提？
(3) \texttt{\#print axioms}／导出文件里的公理是不是只有标准三条？
(4) 证明是否被独立于构建环境的内核重放过？
(5) 被检查的文件与发布的源码是否哈希对应？
五问都过关时，「中间定义看起来很可疑」不再是有效性问题；任何一问不过关，才值得深究。
\item \textbf{给证明作者与维护社区：}删除或显式注释这个未被使用的几何槽（\lean{Parameters.Γ}、\lean{OAt.geo}、\lean{GeoCl}）以及不在依赖锥内的「真几何」代码，
并公布一份「依赖锥内声明清单」。把 1,468,418 个真正被用到的声明与环境中其余约 123 万个常量区分开，是让人类开始理解这份证明的第一步，也能从源头消除此类误读。
\item \textbf{给批评者：}这个观察本身是有价值的代码考古——它准确地指出了形式化证明与其命名所暗示的数学叙事之间的落差。
把它表述为「该证明的几何部分是空壳，实际论证不经过它」是站得住的；把它表述为「该证明无效」则站不住。
\end{itemize}

\appendix
\section{内核级扫描结果明细}
\label{app:conescan}

\subsection{导出文件概况（conescan 第一遍）}

\begin{center}\small
\begin{tabular}{@{}lr@{\qquad}lr@{}}
\toprule
\multicolumn{2}{c}{\textbf{总量}} & \multicolumn{2}{c}{\textbf{声明记录}}\\
\midrule
行数 & 1,080,769,508 & \lean{thm} & 1,246,110 \\
字节数 & 61,570,595,788 & \lean{def} & 214,241 \\
名字 & 14,065,150 & \lean{inductive}（声明块） & 7,969 \\
表达式结点 & 1,065,214,335 & \lean{opaque} & 91 \\
其中 \lean{app} & 989,707,003 & \lean{quot} & 4 \\
其中 \lean{proj} & 65,718 & \lean{axiom} & 3 \\
标记 unsafe／partial 的记录 & 0 & \textbf{合计} & \textbf{1,468,418} \\
\bottomrule
\end{tabular}
\end{center}

公理：\lean{propext}、\lean{Quot.sound}、\lean{Classical.choice}。导出器：lean4export 3.1.0；Lean 4.33.1（githash \texttt{819816b2}）。

\subsection{依赖 \texttt{anyGeometry} 的全部 30 个声明}

「类型／值」两列表示 \lean{anyGeometry}（经传递依赖）出现在该声明的类型里还是证明项（值）里；「直接见证」是该声明直接引用的一个已被污染的常量。

{\scriptsize
\begin{longtable}{@{}rlcc>{\raggedright\arraybackslash}p{6.9cm}>{\raggedright\arraybackslash}p{5.4cm}@{}}
\caption{内核级：传递依赖 \texttt{GlobalTheory.Mem.Top.Bl.anyGeometry} 的声明}\label{tab:tainted}\\
\toprule
\# & 种类 & 类型 & 值 & 声明 & 直接见证 \\
\midrule
\endfirsthead
\toprule
\# & 种类 & 类型 & 值 & 声明 & 直接见证 \\
\midrule
\endhead
\bottomrule
\endfoot
\tableinput{tables/tainted_anygeometry.tex}
\end{longtable}
}

\subsection{读取 \texttt{OAt} 的 \texttt{geo} 字段的声明}

直接读取该字段的只有两个定理：\lean{divEx_geoCl} 与 \lean{oAt_setDiv}（表中「直接见证」为 \lean{OAt.geo} 的两行）；其余各行是它们的传递使用者。

{\scriptsize
\begin{longtable}{@{}rlcc>{\raggedright\arraybackslash}p{6.9cm}>{\raggedright\arraybackslash}p{5.4cm}@{}}
\caption{内核级：传递依赖「对 \texttt{OAt} 第 0 字段（\texttt{geo}）的投影」的声明}\\
\toprule
\# & 种类 & 类型 & 值 & 声明 & 直接见证 \\
\midrule
\endfirsthead
\toprule
\# & 种类 & 类型 & 值 & 声明 & 直接见证 \\
\midrule
\endhead
\bottomrule
\endfoot
\tableinput{tables/tainted_oatgeo.tex}
\end{longtable}
}

\subsection{引用 \texttt{GeoCl.rec} 的声明}

只有 \lean{GeoCl} 自身的归纳声明块（类型、两个构造子与递归子共 4 条记录，递归子的类型提到它自己）。没有任何其他声明。

\scanThreeAppendix

\section{其余源码摘录}

从 \path{Main.lean} 到「文献定理」层的三级转接：

\excerpt{maintheorems}
\excerpt{mainconcl}
\excerpt{lit810}

\section{完整哈希与版本}

\begin{center}\scriptsize
\renewcommand{\arraystretch}{1.25}
\begin{tabularx}{\textwidth}{@{}>{\raggedright\arraybackslash}p{3.6cm}>{\raggedright\arraybackslash}X@{}}
\toprule
对象 & SHA-256／提交 \\
\midrule
\path{lean.tar.gz}（作者原始包） & \texttt{9631b4f020199e580ba86178cb56c6e9b09f690c2369a80dc4334aa93f49c7fc} \\
\path{solution.ndjson.zst} & \texttt{61a8938ad5f5e57397c1c0e3b954916af546d9004ade01c3555744a681b05381} \\
\path{solution.ndjson} & \texttt{5ea74b912147c5c93cc2b6d70cc727cf908df2f62ac6a20480a3fa02f6be244e} \\
\path{challenge.ndjson}（公布的与本机重建的） & \texttt{bd484ac9886a96f26b1d9a7a59ebd45b7a70d102c015d4f31f9302bbdd6ab2b9} \\
镜像 \path{the-omega-institute/riemann-hypothesis} & 提交 \texttt{c7d25b3bfd7906388bb8eb0774e027a2b756a7b6} \\
Mathlib & \texttt{0df444a360eaa60ab8c11dca51a86af692955474}（官方标签 \texttt{v4.33.1}） \\
con-leche & \texttt{bc6676101c488e21ce71a317b33905534d497e14}，以 Lean \texttt{v4.36.0}（\texttt{611c117b}）编译 \\
lean4export & \texttt{15f6055e299ad5b89345e533cc2192f4cc00f659}（标签 \texttt{v4.33.0}），以 Lean \texttt{v4.33.1}（\texttt{819816b2}）编译 \\
\bottomrule
\end{tabularx}
\end{center}

\section{复现步骤}

仓库 \path{evidence/} 下的 \path{REPRODUCE.md} 给出了逐条命令。概要如下（需要约 80\,GB 磁盘；con-leche 一步在 16 核机器上约需一小时）：

\begin{Verbatim}[fontsize=\scriptsize,frame=single,rulecolor=\color{gray!50}]
# 1. 源码：下载 alpo.ge 的 14 个分卷，拼接并核对 SHA-256，解包
for i in 01 02 ... 14; do curl -LO https://alpo.ge/lean.tar.gz.part$i; done
cat lean.tar.gz.part?? > lean.tar.gz && shasum -a 256 lean.tar.gz

# 2. 源码层面：定位定义与调用链
rg -n -w anyGeometry src | wc -l          # 118 行，24 个文件
rg -n 'everyE₂|plugAll' src               # plugAll 只被 everyE₂ 调用；everyE₂ 无调用者

# 3. 内核层面：下载导出文件，核对哈希，运行 conescan
curl -LO https://math.mit.edu/~drew/grh/solution.ndjson.zst
zstd -dc --long=30 solution.ndjson.zst | conescan OUTDIR seeds.tsv watch.txt

# 4. 独立内核重放与命题比对
con-leche --verified --compare challenge.ndjson solution.ndjson

# 5. 独立重建挑战导出（官方 Lean v4.33.1 + Mathlib 官方缓存 + lean4export v4.33.0）
lake exe cache get && lake build Challenge
lake env lean4export Challenge -- $(cat export-targets.txt) | shasum -a 256
\end{Verbatim}

\section{参考来源}

{\small\sloppy
\begin{itemize}
\item L.~Alpöge, 公告：\url{https://alpo.ge/remark.pdf}；源码分卷：\url{https://alpo.ge/lean.tar.gz.part01} 至 \texttt{part14}。
\item J.~L.~Nielsen, \emph{Invalidity of the Alpoge Proposed Proof of the Riemann Hypothesis}：\url{https://philarchive.org/rec/NIEIOT-2}（本次审计未能取得全文）；
作者在 X 上的表述：\url{https://x.com/QualiaQuanta/status/2123871320949825776}。
\item frankchu, \emph{GRH proof's geometry is \texttt{fun \_ \_ \_ => True}, in 24 files}, DEV Community：\url{https://dev.to/frankchu/grh-proofs-geometry-is-fun-true-in-24-files-2n2b}。
\item Lean 与 Mathlib 社区联合声明：\url{https://leanprover-community.github.io/blog/posts/riemann-hypothesis/}。
\item A.~Sutherland, \emph{Verification files for the Lean proof of the Riemann Hypothesis}：\url{https://math.mit.edu/~drew/grh/}。
\item 源码镜像：\url{https://github.com/the-omega-institute/riemann-hypothesis}；社区验证仓库：\url{https://github.com/Deicyde/riemann-hypothesis}。
\item 检查器：\url{https://github.com/leanprover/con-leche}；导出器：\url{https://github.com/leanprover/lean4export}。
\end{itemize}
}

\end{document}
